%% ECON 282E -- Session 5, Deck A5: Higher Order, Risk, and Pruning
%% Source: Quant_Macro Lec.9 frames 31-36 (Higher Order; Extension: Perturbation
%% Around a Steady-State Distribution), plus material absent from that deck:
%% Schmitt-Grohe & Uribe (2004), the odd-derivative result and the risk
%% correction (Lecture_SM_5_Perturbation_I.pdf pp.21, 23, 24), the radius of
%% convergence (pp.31-33), regular vs singular perturbation (p.7), and pruning
%% (Kim-Kim-Schaumburg-Sims 2008; Andreasen-FV-Rubio-Ramirez 2018), which appears
%% nowhere in any source tree.
%% NOTATION: the source's mu (distribution) is lambda here, and its T is Q.
%% Numbers from tools/figures/s05_numbers.json, keys "higher_order", "pruning".

%% ECON 282E -- Foundations of Macroeconomics (UC Riverside, Fall 2026)
%% Shared Beamer preamble for every deck in slides/S01 ... slides/S10.
%%
%% Usage (a deck lives in slides/S0X/ and is compiled from that folder):
%%   %% ECON 282E -- Foundations of Macroeconomics (UC Riverside, Fall 2026)
%% Shared Beamer preamble for every deck in slides/S01 ... slides/S10.
%%
%% Usage (a deck lives in slides/S0X/ and is compiled from that folder):
%%   %% ECON 282E -- Foundations of Macroeconomics (UC Riverside, Fall 2026)
%% Shared Beamer preamble for every deck in slides/S01 ... slides/S10.
%%
%% Usage (a deck lives in slides/S0X/ and is compiled from that folder):
%%   \input{../preamble/course_preamble.tex}
%%   \decktitle{1}{A1}{Who I Am \& Why This Course}{September 24, 2026}
%%   \begin{document}
%%   \makedecktitle
%%   ...
%%
%% Adapted from Courses/AI-ECON-2026/source/shared_preamble.tex (Metropolis theme,
%% concept/caution/example/takeaway boxes, \sectiondivider). Changes for this course:
%%   * course title block (\decktitle / \makedecktitle) instead of \makebeamertitle
%%   * \zh{...} typesets Chinese (book titles) under pdflatex via CJKutf8
%%   * researchbox for the "Research in the wild" frames
%%   * section pages off; TikZ libraries and the course map loaded here
%% Build with pdflatex only (two passes); tools/build_slides.py does this for you.

\documentclass[10pt,english,aspectratio=169]{beamer}

\usepackage[T1]{fontenc}
\usepackage{textcomp}
\usepackage[utf8]{inputenc}
\usepackage{babel}
\usepackage{CJKutf8}

\usepackage{amsmath, amssymb, amsthm, graphicx}
\usepackage{booktabs, multirow, tabularx}
\usepackage{tikz}
\usetikzlibrary{positioning,arrows.meta,shapes,calc,fit,backgrounds}
\usepackage{pgfplots}
\pgfplotsset{compat=1.16}
\usepackage[most]{tcolorbox}
\tcbuselibrary{skins, breakable}
\usepackage{listings}
\usepackage{xcolor}

\lstset{
  basicstyle=\ttfamily\small,
  keywordstyle=\color{blue!80!black}\bfseries,
  commentstyle=\color{green!50!black}\itshape,
  stringstyle=\color{orange!80!black},
  backgroundcolor=\color{gray!8},
  frame=single,
  rulecolor=\color{gray!40},
  breaklines=true,
  showstringspaces=false,
  numbers=none,
}

\usetheme{metropolis}
\usefonttheme{professionalfonts}
\metroset{sectionpage=none}

\definecolor{LightGray}{RGB}{230,230,230}
\definecolor{RedTitle}{RGB}{0,0,200}
\definecolor{VeryLightGray}{RGB}{245,245,245}
\definecolor{DarkText}{RGB}{0,0,0}
\definecolor{UCRBlue}{RGB}{0,45,106}
\definecolor{UCRGold}{RGB}{241,171,0}

% Stage colours of the course arc (used by the course map and elsewhere)
\colorlet{StageTrunk}{blue!12}
\colorlet{StageBranch}{cyan!14}
\colorlet{StageMachine}{gray!18}
\colorlet{StageWall}{red!14}
\colorlet{StageTools}{orange!20}
\colorlet{StageData}{green!16}
\colorlet{StageAgents}{violet!14}

\hypersetup{bookmarks=false,breaklinks=true,pdfborder={0 0 0},colorlinks=true,
  linkcolor=DarkText,urlcolor=RedTitle}

\setbeamercolor{background canvas}{bg=white}
\setbeamercolor{title}{fg=RedTitle}
\setbeamercolor{frametitle}{bg=VeryLightGray,fg=DarkText}
\setbeamercolor{normal text}{fg=DarkText}
\setbeamercolor{alerted text}{fg=RedTitle}
\setbeamersize{text margin left=5mm, text margin right=5mm}
\addtobeamertemplate{frametitle}{}{\vspace{-0.5em}}

\setbeamertemplate{navigation symbols}{}
\setbeamertemplate{footline}{%
  \hspace{5mm}{\usebeamerfont{page number in head/foot}\color{black!45}%
  ECON 282E \,$\cdot$\, \insertshortdeck}\hfill%
  \usebeamercolor[fg]{page number in head/foot}%
  \usebeamerfont{page number in head/foot}%
  \insertframenumber\,/\,\inserttotalframenumber\kern1em\vskip2pt%
}

% ---------------------------------------------------------------------------
% Chinese text under pdflatex (book titles etc.)
\newcommand{\zh}[1]{\begin{CJK*}{UTF8}{gbsn}#1\end{CJK*}}

% ---------------------------------------------------------------------------
% Deck title block
%   \decktitle{session}{deck id}{title}{date}
\newcommand{\insertshortdeck}{}
\newcommand{\decksession}{}
\newcommand{\deckid}{}
\newcommand{\decktitletext}{}
\newcommand{\deckdate}{}
\newcommand{\decktitle}[4]{%
  \renewcommand{\decksession}{#1}%
  \renewcommand{\deckid}{#2}%
  \renewcommand{\decktitletext}{#3}%
  \renewcommand{\deckdate}{#4}%
  \renewcommand{\insertshortdeck}{Session #1 \,$\cdot$\, Deck #1.#2}%
  \title{#3}%
  \author{Zhigang Feng}%
  \date{#4}%
}
\newcommand{\makedecktitle}{%
  \begin{frame}[plain,noframenumbering]
    \begin{tikzpicture}[remember picture,overlay]
      \fill[UCRBlue] (current page.north west) rectangle ([yshift=-0.55cm]current page.north east);
      \fill[UCRGold] ([yshift=-0.55cm]current page.north west) rectangle ([yshift=-0.62cm]current page.north east);
      \node[anchor=west, text=white, font=\small\bfseries] at ([xshift=5mm,yshift=-0.275cm]current page.north west)
        {ECON 282E \,$\cdot$\, Foundations of Macroeconomics};
      \node[anchor=east, text=white, font=\small] at ([xshift=-5mm,yshift=-0.275cm]current page.north east)
        {UC Riverside \,$\cdot$\, Fall 2026};
    \end{tikzpicture}
    \vspace*{1.1cm}
    {\usebeamerfont{title}\color{black!55}\large Session \decksession\ \,$\cdot$\, Deck \decksession.\deckid\par}
    \vspace{0.25cm}
    {\usebeamerfont{title}\color{RedTitle}\LARGE\bfseries \decktitletext\par}
    \vspace{0.7cm}
    {\large Zhigang Feng\par}
    \vspace{0.15cm}
    {\color{black!65}\deckdate\par}
    \vfill
    {\footnotesize\itshape\color{black!55}Build it on paper $\;\to\;$ solve it on the machine $\;\to\;$ push it to the frontier with AI\par}
  \end{frame}}

% ---------------------------------------------------------------------------
% Section divider (plain frame, not counted)
\newcommand{\sectiondivider}[2]{%
\begin{frame}[plain,noframenumbering]
\begin{center}
\vfill
{\Large\textbf{#1}\par}
\vspace{0.35cm}
{\normalsize #2\par}
\vfill
\end{center}
\end{frame}
}

% ---------------------------------------------------------------------------
% Boxes
\newtcolorbox{conceptbox}[1][]{colback=blue!3, colframe=RedTitle!55, boxrule=0.35mm,
  arc=1mm, left=2mm, right=2mm, top=1mm, bottom=1mm, #1}
\newtcolorbox{cautionbox}[1][]{colback=red!3, colframe=red!50!black, boxrule=0.35mm,
  arc=1mm, left=2mm, right=2mm, top=1mm, bottom=1mm, #1}
\newtcolorbox{examplebox}[1][]{colback=green!3, colframe=green!45!black, boxrule=0.35mm,
  arc=1mm, left=2mm, right=2mm, top=1mm, bottom=1mm, #1}
\newtcolorbox{takeawaybox}[1][]{colback=VeryLightGray, colframe=RedTitle!55, boxrule=0.35mm,
  arc=1mm, left=2mm, right=2mm, top=1mm, bottom=1mm, #1}
\newtcolorbox{researchbox}[1][]{enhanced, colback=violet!4, colframe=violet!60!black,
  boxrule=0.35mm, arc=1mm, left=2mm, right=2mm, top=1mm, bottom=1mm,
  fonttitle=\bfseries\small, title={Research in the wild}, #1}

\newcommand{\compacttables}{\renewcommand{\arraystretch}{1.15}}
\newcommand{\src}[1]{{\tiny\color{black!55}#1}}

% ---------------------------------------------------------------------------
\input{../preamble/course_map.tex}

%%   \decktitle{1}{A1}{Who I Am \& Why This Course}{September 24, 2026}
%%   \begin{document}
%%   \makedecktitle
%%   ...
%%
%% Adapted from Courses/AI-ECON-2026/source/shared_preamble.tex (Metropolis theme,
%% concept/caution/example/takeaway boxes, \sectiondivider). Changes for this course:
%%   * course title block (\decktitle / \makedecktitle) instead of \makebeamertitle
%%   * \zh{...} typesets Chinese (book titles) under pdflatex via CJKutf8
%%   * researchbox for the "Research in the wild" frames
%%   * section pages off; TikZ libraries and the course map loaded here
%% Build with pdflatex only (two passes); tools/build_slides.py does this for you.

\documentclass[10pt,english,aspectratio=169]{beamer}

\usepackage[T1]{fontenc}
\usepackage{textcomp}
\usepackage[utf8]{inputenc}
\usepackage{babel}
\usepackage{CJKutf8}

\usepackage{amsmath, amssymb, amsthm, graphicx}
\usepackage{booktabs, multirow, tabularx}
\usepackage{tikz}
\usetikzlibrary{positioning,arrows.meta,shapes,calc,fit,backgrounds}
\usepackage{pgfplots}
\pgfplotsset{compat=1.16}
\usepackage[most]{tcolorbox}
\tcbuselibrary{skins, breakable}
\usepackage{listings}
\usepackage{xcolor}

\lstset{
  basicstyle=\ttfamily\small,
  keywordstyle=\color{blue!80!black}\bfseries,
  commentstyle=\color{green!50!black}\itshape,
  stringstyle=\color{orange!80!black},
  backgroundcolor=\color{gray!8},
  frame=single,
  rulecolor=\color{gray!40},
  breaklines=true,
  showstringspaces=false,
  numbers=none,
}

\usetheme{metropolis}
\usefonttheme{professionalfonts}
\metroset{sectionpage=none}

\definecolor{LightGray}{RGB}{230,230,230}
\definecolor{RedTitle}{RGB}{0,0,200}
\definecolor{VeryLightGray}{RGB}{245,245,245}
\definecolor{DarkText}{RGB}{0,0,0}
\definecolor{UCRBlue}{RGB}{0,45,106}
\definecolor{UCRGold}{RGB}{241,171,0}

% Stage colours of the course arc (used by the course map and elsewhere)
\colorlet{StageTrunk}{blue!12}
\colorlet{StageBranch}{cyan!14}
\colorlet{StageMachine}{gray!18}
\colorlet{StageWall}{red!14}
\colorlet{StageTools}{orange!20}
\colorlet{StageData}{green!16}
\colorlet{StageAgents}{violet!14}

\hypersetup{bookmarks=false,breaklinks=true,pdfborder={0 0 0},colorlinks=true,
  linkcolor=DarkText,urlcolor=RedTitle}

\setbeamercolor{background canvas}{bg=white}
\setbeamercolor{title}{fg=RedTitle}
\setbeamercolor{frametitle}{bg=VeryLightGray,fg=DarkText}
\setbeamercolor{normal text}{fg=DarkText}
\setbeamercolor{alerted text}{fg=RedTitle}
\setbeamersize{text margin left=5mm, text margin right=5mm}
\addtobeamertemplate{frametitle}{}{\vspace{-0.5em}}

\setbeamertemplate{navigation symbols}{}
\setbeamertemplate{footline}{%
  \hspace{5mm}{\usebeamerfont{page number in head/foot}\color{black!45}%
  ECON 282E \,$\cdot$\, \insertshortdeck}\hfill%
  \usebeamercolor[fg]{page number in head/foot}%
  \usebeamerfont{page number in head/foot}%
  \insertframenumber\,/\,\inserttotalframenumber\kern1em\vskip2pt%
}

% ---------------------------------------------------------------------------
% Chinese text under pdflatex (book titles etc.)
\newcommand{\zh}[1]{\begin{CJK*}{UTF8}{gbsn}#1\end{CJK*}}

% ---------------------------------------------------------------------------
% Deck title block
%   \decktitle{session}{deck id}{title}{date}
\newcommand{\insertshortdeck}{}
\newcommand{\decksession}{}
\newcommand{\deckid}{}
\newcommand{\decktitletext}{}
\newcommand{\deckdate}{}
\newcommand{\decktitle}[4]{%
  \renewcommand{\decksession}{#1}%
  \renewcommand{\deckid}{#2}%
  \renewcommand{\decktitletext}{#3}%
  \renewcommand{\deckdate}{#4}%
  \renewcommand{\insertshortdeck}{Session #1 \,$\cdot$\, Deck #1.#2}%
  \title{#3}%
  \author{Zhigang Feng}%
  \date{#4}%
}
\newcommand{\makedecktitle}{%
  \begin{frame}[plain,noframenumbering]
    \begin{tikzpicture}[remember picture,overlay]
      \fill[UCRBlue] (current page.north west) rectangle ([yshift=-0.55cm]current page.north east);
      \fill[UCRGold] ([yshift=-0.55cm]current page.north west) rectangle ([yshift=-0.62cm]current page.north east);
      \node[anchor=west, text=white, font=\small\bfseries] at ([xshift=5mm,yshift=-0.275cm]current page.north west)
        {ECON 282E \,$\cdot$\, Foundations of Macroeconomics};
      \node[anchor=east, text=white, font=\small] at ([xshift=-5mm,yshift=-0.275cm]current page.north east)
        {UC Riverside \,$\cdot$\, Fall 2026};
    \end{tikzpicture}
    \vspace*{1.1cm}
    {\usebeamerfont{title}\color{black!55}\large Session \decksession\ \,$\cdot$\, Deck \decksession.\deckid\par}
    \vspace{0.25cm}
    {\usebeamerfont{title}\color{RedTitle}\LARGE\bfseries \decktitletext\par}
    \vspace{0.7cm}
    {\large Zhigang Feng\par}
    \vspace{0.15cm}
    {\color{black!65}\deckdate\par}
    \vfill
    {\footnotesize\itshape\color{black!55}Build it on paper $\;\to\;$ solve it on the machine $\;\to\;$ push it to the frontier with AI\par}
  \end{frame}}

% ---------------------------------------------------------------------------
% Section divider (plain frame, not counted)
\newcommand{\sectiondivider}[2]{%
\begin{frame}[plain,noframenumbering]
\begin{center}
\vfill
{\Large\textbf{#1}\par}
\vspace{0.35cm}
{\normalsize #2\par}
\vfill
\end{center}
\end{frame}
}

% ---------------------------------------------------------------------------
% Boxes
\newtcolorbox{conceptbox}[1][]{colback=blue!3, colframe=RedTitle!55, boxrule=0.35mm,
  arc=1mm, left=2mm, right=2mm, top=1mm, bottom=1mm, #1}
\newtcolorbox{cautionbox}[1][]{colback=red!3, colframe=red!50!black, boxrule=0.35mm,
  arc=1mm, left=2mm, right=2mm, top=1mm, bottom=1mm, #1}
\newtcolorbox{examplebox}[1][]{colback=green!3, colframe=green!45!black, boxrule=0.35mm,
  arc=1mm, left=2mm, right=2mm, top=1mm, bottom=1mm, #1}
\newtcolorbox{takeawaybox}[1][]{colback=VeryLightGray, colframe=RedTitle!55, boxrule=0.35mm,
  arc=1mm, left=2mm, right=2mm, top=1mm, bottom=1mm, #1}
\newtcolorbox{researchbox}[1][]{enhanced, colback=violet!4, colframe=violet!60!black,
  boxrule=0.35mm, arc=1mm, left=2mm, right=2mm, top=1mm, bottom=1mm,
  fonttitle=\bfseries\small, title={Research in the wild}, #1}

\newcommand{\compacttables}{\renewcommand{\arraystretch}{1.15}}
\newcommand{\src}[1]{{\tiny\color{black!55}#1}}

% ---------------------------------------------------------------------------
%% Course map: the recurring "you are here" slide for ECON 282E.
%% Loaded by course_preamble.tex. Pattern follows AI-ECON-2026 lec01_map.tex, but the map
%% shows the whole ten-session course rather than one lecture.
%%
%%   \CourseMap{s}{label}   s = 1..10 highlights session s and prints "you are here: label"
%%                          (e.g. \CourseMap{1}{1.A2}); earlier sessions get a check mark.
%%   \CourseMap{0}{}        full overview, nothing highlighted.
%%
%% Note: TikZ option lists are comma-split before TeX conditionals run, so every \ifnum
%% below sits outside [...] option lists.

\newcommand{\CourseMap}[2]{%
\begin{center}
\begin{tikzpicture}[
    sess/.style={rectangle, rounded corners=2pt, draw=black!45, minimum width=1.34cm,
        minimum height=1.12cm, text width=1.26cm, align=center, inner sep=1pt},
    stage/.style={font=\tiny\bfseries, text=black!70},
]
  \foreach \i/\d/\t/\c in {%
      1/Sep 24/Intro \\ Growth/StageTrunk,
      2/Oct 1/HA $\cdot$ OLG \\ Policy/StageBranch,
      3/Oct 8/Num.\ DP \\ Python/StageMachine,
      4/Oct 15/Numerics \\ I/StageMachine,
      5/Oct 22/Numerics \\ II/StageMachine,
      6/Oct 29/The wall \\ \textbf{Midterm}/StageWall,
      7/Nov 5/ML \\ Deep nets/StageTools,
      8/Nov 12/RL \\ HA $+$ DL/StageTools,
      9/Nov 19/Text \\ LLMs/StageData,
      10/Dec 3/Agents \\ \textbf{Projects}/StageAgents} {
    \node[sess, fill=\c] (n\i) at ({1.46*(\i-1)}, 0)
      {{\scriptsize\bfseries S\i}\\[-1pt]{\tiny\color{black!60}\d}\\[1pt]{\tiny \t}};
    \ifnum\i<#1
      \node[font=\scriptsize, text=green!45!black] at ([xshift=-2.5mm,yshift=-2.2mm]n\i.north east) {$\checkmark$};
    \fi
    \ifnum\i=#1
      \draw[RedTitle, very thick, rounded corners=2pt]
        ([xshift=-0.6mm,yshift=-0.6mm]n\i.south west) rectangle ([xshift=0.6mm,yshift=0.6mm]n\i.north east);
      % keep the label inside the picture at both ends of the row
      \ifnum\i<3
        \node[font=\scriptsize\bfseries, text=RedTitle, anchor=south west] at ([yshift=1.2mm]n\i.north west)
          {$\blacktriangledown$\,you are here: #2};
      \else\ifnum\i>8
        \node[font=\scriptsize\bfseries, text=RedTitle, anchor=south east] at ([yshift=1.2mm]n\i.north east)
          {you are here: #2\,$\blacktriangledown$};
      \else
        \node[font=\scriptsize\bfseries, text=RedTitle, anchor=south] at ([yshift=1.2mm]n\i.north)
          {you are here: #2\,$\blacktriangledown$};
      \fi\fi
    \fi
  }
  % stage band under the sessions
  \foreach \a/\b/\lab/\c in {1/1/trunk/StageTrunk, 2/2/branches/StageBranch,
      3/5/the machine $\cdot$ classical methods/StageMachine, 6/6/the wall/StageWall,
      7/8/new tools/StageTools, 9/9/new data/StageData, 10/10/colleagues/StageAgents} {
    \fill[\c] ([yshift=-1.5mm]n\a.south west) rectangle ([yshift=-4.6mm]n\b.south east);
    \path ([yshift=-1.5mm]n\a.south west) -- ([yshift=-4.6mm]n\b.south east)
      node[midway, stage] {\lab};
  }
  \node[font=\footnotesize\itshape, text=black!70] at ([yshift=-9.5mm]$(n1.south)!0.5!(n10.south)$)
    {Build it on paper $\;\to\;$ solve it on the machine $\;\to\;$ push it to the frontier with AI};
\end{tikzpicture}
\end{center}
}


%%   \decktitle{1}{A1}{Who I Am \& Why This Course}{September 24, 2026}
%%   \begin{document}
%%   \makedecktitle
%%   ...
%%
%% Adapted from Courses/AI-ECON-2026/source/shared_preamble.tex (Metropolis theme,
%% concept/caution/example/takeaway boxes, \sectiondivider). Changes for this course:
%%   * course title block (\decktitle / \makedecktitle) instead of \makebeamertitle
%%   * \zh{...} typesets Chinese (book titles) under pdflatex via CJKutf8
%%   * researchbox for the "Research in the wild" frames
%%   * section pages off; TikZ libraries and the course map loaded here
%% Build with pdflatex only (two passes); tools/build_slides.py does this for you.

\documentclass[10pt,english,aspectratio=169]{beamer}

\usepackage[T1]{fontenc}
\usepackage{textcomp}
\usepackage[utf8]{inputenc}
\usepackage{babel}
\usepackage{CJKutf8}

\usepackage{amsmath, amssymb, amsthm, graphicx}
\usepackage{booktabs, multirow, tabularx}
\usepackage{tikz}
\usetikzlibrary{positioning,arrows.meta,shapes,calc,fit,backgrounds}
\usepackage{pgfplots}
\pgfplotsset{compat=1.16}
\usepackage[most]{tcolorbox}
\tcbuselibrary{skins, breakable}
\usepackage{listings}
\usepackage{xcolor}

\lstset{
  basicstyle=\ttfamily\small,
  keywordstyle=\color{blue!80!black}\bfseries,
  commentstyle=\color{green!50!black}\itshape,
  stringstyle=\color{orange!80!black},
  backgroundcolor=\color{gray!8},
  frame=single,
  rulecolor=\color{gray!40},
  breaklines=true,
  showstringspaces=false,
  numbers=none,
}

\usetheme{metropolis}
\usefonttheme{professionalfonts}
\metroset{sectionpage=none}

\definecolor{LightGray}{RGB}{230,230,230}
\definecolor{RedTitle}{RGB}{0,0,200}
\definecolor{VeryLightGray}{RGB}{245,245,245}
\definecolor{DarkText}{RGB}{0,0,0}
\definecolor{UCRBlue}{RGB}{0,45,106}
\definecolor{UCRGold}{RGB}{241,171,0}

% Stage colours of the course arc (used by the course map and elsewhere)
\colorlet{StageTrunk}{blue!12}
\colorlet{StageBranch}{cyan!14}
\colorlet{StageMachine}{gray!18}
\colorlet{StageWall}{red!14}
\colorlet{StageTools}{orange!20}
\colorlet{StageData}{green!16}
\colorlet{StageAgents}{violet!14}

\hypersetup{bookmarks=false,breaklinks=true,pdfborder={0 0 0},colorlinks=true,
  linkcolor=DarkText,urlcolor=RedTitle}

\setbeamercolor{background canvas}{bg=white}
\setbeamercolor{title}{fg=RedTitle}
\setbeamercolor{frametitle}{bg=VeryLightGray,fg=DarkText}
\setbeamercolor{normal text}{fg=DarkText}
\setbeamercolor{alerted text}{fg=RedTitle}
\setbeamersize{text margin left=5mm, text margin right=5mm}
\addtobeamertemplate{frametitle}{}{\vspace{-0.5em}}

\setbeamertemplate{navigation symbols}{}
\setbeamertemplate{footline}{%
  \hspace{5mm}{\usebeamerfont{page number in head/foot}\color{black!45}%
  ECON 282E \,$\cdot$\, \insertshortdeck}\hfill%
  \usebeamercolor[fg]{page number in head/foot}%
  \usebeamerfont{page number in head/foot}%
  \insertframenumber\,/\,\inserttotalframenumber\kern1em\vskip2pt%
}

% ---------------------------------------------------------------------------
% Chinese text under pdflatex (book titles etc.)
\newcommand{\zh}[1]{\begin{CJK*}{UTF8}{gbsn}#1\end{CJK*}}

% ---------------------------------------------------------------------------
% Deck title block
%   \decktitle{session}{deck id}{title}{date}
\newcommand{\insertshortdeck}{}
\newcommand{\decksession}{}
\newcommand{\deckid}{}
\newcommand{\decktitletext}{}
\newcommand{\deckdate}{}
\newcommand{\decktitle}[4]{%
  \renewcommand{\decksession}{#1}%
  \renewcommand{\deckid}{#2}%
  \renewcommand{\decktitletext}{#3}%
  \renewcommand{\deckdate}{#4}%
  \renewcommand{\insertshortdeck}{Session #1 \,$\cdot$\, Deck #1.#2}%
  \title{#3}%
  \author{Zhigang Feng}%
  \date{#4}%
}
\newcommand{\makedecktitle}{%
  \begin{frame}[plain,noframenumbering]
    \begin{tikzpicture}[remember picture,overlay]
      \fill[UCRBlue] (current page.north west) rectangle ([yshift=-0.55cm]current page.north east);
      \fill[UCRGold] ([yshift=-0.55cm]current page.north west) rectangle ([yshift=-0.62cm]current page.north east);
      \node[anchor=west, text=white, font=\small\bfseries] at ([xshift=5mm,yshift=-0.275cm]current page.north west)
        {ECON 282E \,$\cdot$\, Foundations of Macroeconomics};
      \node[anchor=east, text=white, font=\small] at ([xshift=-5mm,yshift=-0.275cm]current page.north east)
        {UC Riverside \,$\cdot$\, Fall 2026};
    \end{tikzpicture}
    \vspace*{1.1cm}
    {\usebeamerfont{title}\color{black!55}\large Session \decksession\ \,$\cdot$\, Deck \decksession.\deckid\par}
    \vspace{0.25cm}
    {\usebeamerfont{title}\color{RedTitle}\LARGE\bfseries \decktitletext\par}
    \vspace{0.7cm}
    {\large Zhigang Feng\par}
    \vspace{0.15cm}
    {\color{black!65}\deckdate\par}
    \vfill
    {\footnotesize\itshape\color{black!55}Build it on paper $\;\to\;$ solve it on the machine $\;\to\;$ push it to the frontier with AI\par}
  \end{frame}}

% ---------------------------------------------------------------------------
% Section divider (plain frame, not counted)
\newcommand{\sectiondivider}[2]{%
\begin{frame}[plain,noframenumbering]
\begin{center}
\vfill
{\Large\textbf{#1}\par}
\vspace{0.35cm}
{\normalsize #2\par}
\vfill
\end{center}
\end{frame}
}

% ---------------------------------------------------------------------------
% Boxes
\newtcolorbox{conceptbox}[1][]{colback=blue!3, colframe=RedTitle!55, boxrule=0.35mm,
  arc=1mm, left=2mm, right=2mm, top=1mm, bottom=1mm, #1}
\newtcolorbox{cautionbox}[1][]{colback=red!3, colframe=red!50!black, boxrule=0.35mm,
  arc=1mm, left=2mm, right=2mm, top=1mm, bottom=1mm, #1}
\newtcolorbox{examplebox}[1][]{colback=green!3, colframe=green!45!black, boxrule=0.35mm,
  arc=1mm, left=2mm, right=2mm, top=1mm, bottom=1mm, #1}
\newtcolorbox{takeawaybox}[1][]{colback=VeryLightGray, colframe=RedTitle!55, boxrule=0.35mm,
  arc=1mm, left=2mm, right=2mm, top=1mm, bottom=1mm, #1}
\newtcolorbox{researchbox}[1][]{enhanced, colback=violet!4, colframe=violet!60!black,
  boxrule=0.35mm, arc=1mm, left=2mm, right=2mm, top=1mm, bottom=1mm,
  fonttitle=\bfseries\small, title={Research in the wild}, #1}

\newcommand{\compacttables}{\renewcommand{\arraystretch}{1.15}}
\newcommand{\src}[1]{{\tiny\color{black!55}#1}}

% ---------------------------------------------------------------------------
%% Course map: the recurring "you are here" slide for ECON 282E.
%% Loaded by course_preamble.tex. Pattern follows AI-ECON-2026 lec01_map.tex, but the map
%% shows the whole ten-session course rather than one lecture.
%%
%%   \CourseMap{s}{label}   s = 1..10 highlights session s and prints "you are here: label"
%%                          (e.g. \CourseMap{1}{1.A2}); earlier sessions get a check mark.
%%   \CourseMap{0}{}        full overview, nothing highlighted.
%%
%% Note: TikZ option lists are comma-split before TeX conditionals run, so every \ifnum
%% below sits outside [...] option lists.

\newcommand{\CourseMap}[2]{%
\begin{center}
\begin{tikzpicture}[
    sess/.style={rectangle, rounded corners=2pt, draw=black!45, minimum width=1.34cm,
        minimum height=1.12cm, text width=1.26cm, align=center, inner sep=1pt},
    stage/.style={font=\tiny\bfseries, text=black!70},
]
  \foreach \i/\d/\t/\c in {%
      1/Sep 24/Intro \\ Growth/StageTrunk,
      2/Oct 1/HA $\cdot$ OLG \\ Policy/StageBranch,
      3/Oct 8/Num.\ DP \\ Python/StageMachine,
      4/Oct 15/Numerics \\ I/StageMachine,
      5/Oct 22/Numerics \\ II/StageMachine,
      6/Oct 29/The wall \\ \textbf{Midterm}/StageWall,
      7/Nov 5/ML \\ Deep nets/StageTools,
      8/Nov 12/RL \\ HA $+$ DL/StageTools,
      9/Nov 19/Text \\ LLMs/StageData,
      10/Dec 3/Agents \\ \textbf{Projects}/StageAgents} {
    \node[sess, fill=\c] (n\i) at ({1.46*(\i-1)}, 0)
      {{\scriptsize\bfseries S\i}\\[-1pt]{\tiny\color{black!60}\d}\\[1pt]{\tiny \t}};
    \ifnum\i<#1
      \node[font=\scriptsize, text=green!45!black] at ([xshift=-2.5mm,yshift=-2.2mm]n\i.north east) {$\checkmark$};
    \fi
    \ifnum\i=#1
      \draw[RedTitle, very thick, rounded corners=2pt]
        ([xshift=-0.6mm,yshift=-0.6mm]n\i.south west) rectangle ([xshift=0.6mm,yshift=0.6mm]n\i.north east);
      % keep the label inside the picture at both ends of the row
      \ifnum\i<3
        \node[font=\scriptsize\bfseries, text=RedTitle, anchor=south west] at ([yshift=1.2mm]n\i.north west)
          {$\blacktriangledown$\,you are here: #2};
      \else\ifnum\i>8
        \node[font=\scriptsize\bfseries, text=RedTitle, anchor=south east] at ([yshift=1.2mm]n\i.north east)
          {you are here: #2\,$\blacktriangledown$};
      \else
        \node[font=\scriptsize\bfseries, text=RedTitle, anchor=south] at ([yshift=1.2mm]n\i.north)
          {you are here: #2\,$\blacktriangledown$};
      \fi\fi
    \fi
  }
  % stage band under the sessions
  \foreach \a/\b/\lab/\c in {1/1/trunk/StageTrunk, 2/2/branches/StageBranch,
      3/5/the machine $\cdot$ classical methods/StageMachine, 6/6/the wall/StageWall,
      7/8/new tools/StageTools, 9/9/new data/StageData, 10/10/colleagues/StageAgents} {
    \fill[\c] ([yshift=-1.5mm]n\a.south west) rectangle ([yshift=-4.6mm]n\b.south east);
    \path ([yshift=-1.5mm]n\a.south west) -- ([yshift=-4.6mm]n\b.south east)
      node[midway, stage] {\lab};
  }
  \node[font=\footnotesize\itshape, text=black!70] at ([yshift=-9.5mm]$(n1.south)!0.5!(n10.south)$)
    {Build it on paper $\;\to\;$ solve it on the machine $\;\to\;$ push it to the frontier with AI};
\end{tikzpicture}
\end{center}
}


\graphicspath{{./figures/}}

\lstset{language=Python, basicstyle=\ttfamily\scriptsize, frame=none,
        backgroundcolor=\color{gray!8}, aboveskip=2pt, belowskip=2pt}

\decktitle{5}{A5}{Higher Order, Risk, and Pruning}{Thursday, October 22, 2026}

\begin{document}

\makedecktitle

%=============================================================================
\begin{frame}{Where We Are}
\CourseMap{5}{5.A5}
\vspace{-1mm}
\begin{center}
\small \textbf{Block A:} 5.A1 what perturbation computes $\;\cdot\;$ 5.A2 the model and the parameter $\;\cdot\;$ 5.A3 first order $\;\cdot\;$ 5.A4 linear RE systems $\;\cdot\;$ \textbf{5.A5} higher order, risk, pruning $\;\cdot\;$ 5.A6 the worked example.
\end{center}
\end{frame}

%=============================================================================
\begin{frame}{Why Go Beyond First Order}
\begin{columns}[T]
\begin{column}{0.54\textwidth}
\small
The first-order solution is \textbf{certainty equivalent}: $c_\chi=k_\chi=0$, so the linear rules are identical in a risky economy and a riskless one. Whatever a linearized model says about uncertainty, it was assumed, not computed.
\medskip

Second and higher order restore it:
\begin{itemize}\small
  \item \textbf{precautionary behaviour}, through $c_{\chi\chi}$;
  \item \textbf{risk premia} and the welfare cost of uncertainty;
  \item \textbf{skewness} effects, at third order.
\end{itemize}
\medskip
The cost is algebra --- and the method is still local, so a higher-order expansion of a badly-behaved model is a more elaborate wrong answer, not a right one.
\end{column}
\begin{column}{0.43\textwidth}
\begin{takeawaybox}\footnotesize
Order buys \textbf{curvature}, not \textbf{range}. If the problem is that you are far from the steady state, going to third order is the wrong tool --- that is what Block B is for.
\end{takeawaybox}
\medskip
\begin{examplebox}\footnotesize
Reference: Schmitt-Groh\'e \& Uribe (2004) for the general second-order solution; Fern\'andez-Villaverde, Rubio-Ram\'irez \& Santos (2006) on what the approximation error does to likelihoods.
\end{examplebox}
\end{column}
\end{columns}
\vspace{1mm}
\src{Quant\_Macro Lec.~9, ``Why Go Beyond First Order?''.}
\end{frame}

%=============================================================================
\begin{frame}{Second Order: Twelve Equations, Twelve Unknowns}
\begin{columns}[T]
\begin{column}{0.5\textwidth}
\small
Take second derivatives of $F(k_t,z_t;\chi)$ at $(k^{*},0;0)$ and set each to zero:
\[
F_{kk}=F_{kz}=F_{k\chi}=F_{zz}=F_{z\chi}=F_{\chi\chi}=0 .
\]
Six conditions, each a 2-vector, so \textbf{12 scalar equations}. The unknowns are the six second-order coefficients of $c$ and the six of $k$ --- \textbf{12}, by Young's theorem, since $c_{kz}=c_{zk}$.
\medskip

\textbf{And the system is linear.} The first-order coefficients are already known and enter as constants; the new unknowns appear only linearly. Every order above the first has this property, which is why third and fourth order are tedious rather than hard.
\end{column}
\begin{column}{0.47\textwidth}
\begin{conceptbox}[title=\textbf{Two things vanish for free}]\small
The cross terms $c_{k\chi}$ and $c_{z\chi}$ are \textbf{zero}, and so are their $k$ counterparts.
\medskip
More generally, \textbf{every odd derivative in $\chi$ is zero}. Risk enters only through even powers of the shock scale --- which is to say, through variances.
\end{conceptbox}
\end{column}
\end{columns}
\vspace{1mm}
\src{Quant\_Macro Lec.~9, ``Second-order approximation''; the odd-derivative result and Young's-theorem count from FV\&G, \emph{Solution Methods} V, ``Second-order approximation''.}
\end{frame}

%=============================================================================
\begin{frame}{The Correction for Risk}
\begin{columns}[T]
\begin{column}{0.54\textwidth}
\small
The surviving new object is
\[
  \tfrac12\,c_{\chi\chi}(k^{*},0;0)\,\chi^{2} ,
\]
a \textbf{constant} added to the policy function. At $\chi=1$ it shifts consumption away from its deterministic level for no reason other than the presence of risk.
\medskip

Three consequences:
\begin{itemize}\small
  \item \textbf{Certainty equivalence is gone.} The policy now depends on $\sigma$.
  \item The \textbf{ergodic distribution} of the states moves: the economy no longer fluctuates around the deterministic steady state but around a risk-adjusted one.
  \item Welfare comparisons between economies with different volatility become meaningful for the first time.
\end{itemize}
\end{column}
\begin{column}{0.43\textwidth}
\begin{cautionbox}[title=\textbf{A distinction worth keeping}]\small
\textbf{Risk aversion} is about $u''<0$ --- the second derivative. \textbf{Precautionary saving} is about $u'''>0$ --- the third.
\medskip
A model can have a very risk-averse agent and no precautionary motive at all. Leland (1968), Kimball (1990); and 2.A1, where the same distinction decided who is constrained.
\end{cautionbox}
\end{column}
\end{columns}
\vspace{1mm}
\src{FV\&G, \emph{Solution Methods} V, ``Correction for risk'' and ``Solving the system III''.}
\end{frame}

%=============================================================================
\begin{frame}{Measured: The Second-Order Solution}
\begin{columns}[T]
\begin{column}{0.5\textwidth}
\small
The twelve conditions of the previous frame, imposed on the quarterly model. The six numbers below
are the \textbf{curvature} of the policy functions --- how their slopes change as the state moves.
Worst violation of any of the twelve conditions: $3.6\times10^{-9}$.
\medskip
\begin{center}\footnotesize
\begin{tabularx}{\linewidth}{@{}>{\raggedright\arraybackslash}X r r@{}}
\toprule
 & $\hat{c}$ & $\hat{k}'$ \\
\midrule
on $\hat{k}^{2}$ & $0.0491$ & $0.0287$ \\
on $\hat{k}z$ & $-0.0890$ & $-0.0502$ \\
on $z^{2}$ & $0.1082$ & $0.0827$ \\
\midrule
constant, \% of s.s. & $\mathbf{-0.00041}$ & $+0.000033$ \\
\bottomrule
\end{tabularx}
\end{center}
\medskip
The constant is the risk correction. Consumption is \textbf{lower} than the deterministic solution says, and capital higher --- precautionary saving, with a number on it.
\end{column}
\begin{column}{0.47\textwidth}
\begin{takeawaybox}\footnotesize
Adding this entire block moves the four \emph{first-order} coefficients by $1.5\times10^{-5}$. Certainty equivalence is not an approximation here; the linear terms really are unaffected.
\end{takeawaybox}
\medskip
\begin{examplebox}\footnotesize
$-0.00041\%$ is tiny because $\sigma=0.007$ is tiny. The magnitude is a fact about the calibration; the \emph{sign} and the mechanism are facts about the model.
\end{examplebox}
\end{column}
\end{columns}
\vspace{1mm}
\src{Ledger \texttt{higher\_order}.}
\end{frame}

%=============================================================================
\begin{frame}{Measured: The Correction Scales Like $\sigma^{2}$ --- Until It Doesn't}
\begin{columns}[T]
\begin{column}{0.5\textwidth}
\small
If risk enters only through even powers of $\chi$, the leading term should be quadratic in the shock
size. Re-solving at multiples of $\sigma$, against the prediction ``baseline $\times$
(multiple)$^{2}$'':
\medskip
\begin{center}\footnotesize
\begin{tabularx}{\linewidth}{@{}r r r@{}}
\toprule
$\sigma$ & \textbf{measured} & $-0.00041\times m^{2}$ \\
\midrule
$1\times$  & $-0.00041\%$ & --- \\
$5\times$  & $-0.00998\%$ & $-0.01022\%$ \\
$10\times$ & $-0.03698\%$ & $-0.04088\%$ \\
$20\times$ & $-0.10107\%$ & $-0.16351\%$ \\
\bottomrule
\end{tabularx}
\end{center}
\medskip
Both columns are \% of steady-state consumption. At five times the shock the quadratic law is good
to $2\%$; at twenty times it over-predicts by $60\%$.
\end{column}
\begin{column}{0.47\textwidth}
\begin{conceptbox}[title=\textbf{Read the failure, not the success}]\small
The $\sigma^{2}$ law is a statement about the \emph{leading} term. When the shock is large enough, the fourth-order term the expansion dropped is no longer negligible --- and the second-order solution is the thing being wrong.
\end{conceptbox}
\medskip
\begin{cautionbox}\footnotesize
Extrapolating a perturbation result to a larger shock is exactly the move the method does not license.
\end{cautionbox}
\end{column}
\end{columns}
\vspace{1mm}
\src{Ledger \texttt{higher\_order.risk\_sweep}; all four rows solved to residual $\le 3.6\times10^{-9}$.}
\end{frame}

%=============================================================================
\begin{frame}{Why Second-Order Simulations Explode}
\begin{columns}[T]
\begin{column}{0.55\textwidth}
\small
The second-order policy is
\[
  \hat{k}_{t+1}= b_0 + b_1\hat{k}_t + b_2 z_t
  + \tfrac12\big(b_3\hat{k}_t^{2}+2b_4\hat{k}_tz_t+b_5z_t^{2}\big).
\]
Simulate it by feeding $\hat{k}_{t+1}$ back in as $\hat{k}_t$. That squares a quantity that already contains a squared term, so the recursion generates terms of order $4$, $8$, $16,\dots$ in the shock --- \textbf{orders the approximation never claimed to be valid at}.
\medskip

Sooner or later a sequence of draws pushes the state where the spurious high-order terms dominate, and the path runs away. The explosion is an artifact of the simulation, not a property of the model.
\end{column}
\begin{column}{0.42\textwidth}
\begin{cautionbox}[title=\textbf{Not a stability problem}]\small
$|b_1|<1$: the linear part is stable. It is the \emph{quadratic feedback} that diverges, so no amount of checking eigenvalues warns you.
\end{cautionbox}
\medskip
\begin{examplebox}\footnotesize
It also contaminates moments long before it explodes: a simulation that has not yet diverged can still have a badly wrong variance.
\end{examplebox}
\end{column}
\end{columns}
\end{frame}

%=============================================================================
\begin{frame}{Pruning}
\begin{columns}[T]
\begin{column}{0.55\textwidth}
\small
The fix of \textbf{Kim, Kim, Schaumburg \& Sims (2008)}: never let the simulation carry a term above the order you solved for.
\medskip

Keep \emph{two} state variables. A first-order state,
\[
  \hat{k}^{(1)}_{t+1}=b_1\hat{k}^{(1)}_t+b_2z_t ,
\]
and a second-order correction that is fed \textbf{only first-order terms}:
\[
  \hat{k}^{(2)}_{t+1}=b_0+b_1\hat{k}^{(2)}_t
  +\tfrac12\big(b_3(\hat{k}^{(1)}_t)^{2}+2b_4\hat{k}^{(1)}_tz_t+b_5z_t^{2}\big).
\]
Report $\hat{k}_t=\hat{k}^{(1)}_t+\hat{k}^{(2)}_t$. Nothing is ever squared twice, so no spurious order is ever created.
\end{column}
\begin{column}{0.42\textwidth}
\begin{conceptbox}[title=\textbf{What it is really doing}]\small
Making the simulation consistent with the \textbf{order of the approximation}. The unpruned recursion silently claims accuracy the solution never had.
\end{conceptbox}
\medskip
\begin{examplebox}\footnotesize
Andreasen, Fern\'andez-Villaverde \& Rubio-Ram\'irez (2018) give the general pruned state-space form at any order, with closed-form moments and a likelihood.
\end{examplebox}
\end{column}
\end{columns}
\end{frame}

%=============================================================================
\begin{frame}{Measured: When Pruning Actually Matters}
\begin{columns}[T]
\begin{column}{0.5\textwidth}
\small
$20{,}000$ periods, same shocks, pruned and unpruned:
\medskip
\begin{center}\footnotesize
\begin{tabularx}{\linewidth}{@{}r >{\raggedright\arraybackslash}X >{\raggedright\arraybackslash}X@{}}
\toprule
$\sigma$ & \textbf{unpruned} $\max|\hat{k}|$ & \textbf{pruned} \\
\midrule
$1\times$  & $0.1013$ & $0.1013$ \\
$10\times$ & $0.9947$ & $0.9962$ \\
$30\times$ & \textbf{exploded, $t=3{,}967$} & $3.25$ \\
$60\times$ & \textbf{exploded, $t=166$} & $8.24$ \\
\bottomrule
\end{tabularx}
\end{center}
\medskip
\textbf{At this model's own calibration, pruning changes nothing} --- four significant figures agree. You need thirty times the shock before the unpruned recursion blows up at all.
\end{column}
\begin{column}{0.47\textwidth}
\begin{takeawaybox}\footnotesize
So prune anyway. It costs one extra state variable, it never hurts, and the failure it prevents is \textbf{silent until it is catastrophic} --- and the threshold depends on the model, not on a rule of thumb you can carry between papers.
\end{takeawaybox}
\medskip
\begin{cautionbox}\footnotesize
``It ran without exploding'' is not evidence the moments are right.
\end{cautionbox}
\end{column}
\end{columns}
\vspace{1mm}
\src{Ledger \texttt{pruning}; second-order solves at every $\sigma$ to residual $\le 7.1\times10^{-9}$.}
\end{frame}

%=============================================================================
\begin{frame}{How Far Away Is Any of This Valid?}
\begin{columns}[T]
\begin{column}{0.55\textwidth}
\small
Perturbation is valid within the power series' \textbf{radius of convergence}. There is a clean answer to how big that is, and it is not a real-analysis answer.
\medskip

\begin{conceptbox}[title=\textbf{From complex analysis}]\small
The radius of convergence equals the distance from the expansion point to the \textbf{nearest singularity of the decision rule in the complex plane} --- a pole, a branch point, a discontinuity in the function or any of its derivatives. If there is none, the radius is infinite.
\end{conceptbox}
\medskip
So the trustworthy range is set by a feature of the policy function that may sit nowhere near the real line, and that no amount of staring at a simulation will reveal.
\end{column}
\begin{column}{0.42\textwidth}
\begin{takeawaybox}\footnotesize
Which is why the honest procedure is the one 3.A1 established: \textbf{compute Euler residuals off the expansion point} and look. 5.A6 does exactly that.
\end{takeawaybox}
\medskip
\begin{cautionbox}\footnotesize
And the usual killer is a \textbf{boundary}: a borrowing limit or an irreversibility constraint is a non-removable singularity sitting right where the interesting behaviour is.
\end{cautionbox}
\end{column}
\end{columns}
\vspace{1mm}
\src{FV\&G, \emph{Solution Methods} V, ``Local properties of the solution I--III''.}
\end{frame}

%=============================================================================
\begin{frame}{Regular and Singular Perturbations}
\begin{columns}[T]
\begin{column}{0.55\textwidth}
\small
There is a prior question, and it is not a technicality.
\medskip

\textbf{Regular perturbation:} a small change in the problem produces a small change in the solution. Everything above assumes this.
\medskip

\textbf{Singular perturbation:} a small change in the problem produces a \textbf{large} change in the solution. Then no expansion in the perturbation parameter can work, because the object being expanded is not continuous in it at zero.
\medskip

Most problems in economics are regular. The standard counterexample is squarely in this course's territory: \textbf{introducing a new asset into an incomplete-market model}. Adding an asset with infinitesimal supply can change the equilibrium allocation discontinuously.
\end{column}
\begin{column}{0.42\textwidth}
\begin{cautionbox}[title=\textbf{Why this matters here}]\small
2.A1's whole point was that market structure is an \emph{assumption}. This says the assumption cannot always be perturbed away --- so ``nearly complete markets'' is not always near complete markets.
\end{cautionbox}
\medskip
\begin{examplebox}\footnotesize
Nothing warns you. A singular perturbation produces numbers, and they look like numbers.
\end{examplebox}
\end{column}
\end{columns}
\vspace{1mm}
\src{FV\&G, \emph{Solution Methods} V, ``Regular versus singular perturbations''.}
\end{frame}

%=============================================================================
\begin{frame}{From a Point to a Distribution}
\begin{columns}[T]
\begin{column}{0.55\textwidth}
\small
Everything so far expands around a steady-state \textbf{point}: $(k^{*},c^{*},0)$, one number per variable.
\medskip

In a heterogeneous-agent economy the aggregate state includes the \textbf{cross-sectional distribution} $\lambda$ over wealth --- 2.A3's object. The expansion point is then a steady-state \emph{distribution} $\bar{\lambda}$:
\[
  \lambda_t \;\approx\; \bar{\lambda} \;+\; \text{(first-order deviation)} .
\]
\medskip
The logic does not change at all. Taylor expand around a reference point; the reference point is simply no longer a scalar.
\end{column}
\begin{column}{0.42\textwidth}
\begin{center}\footnotesize
\begin{tabularx}{\linewidth}{@{}>{\raggedright\arraybackslash}p{1.5cm} >{\raggedright\arraybackslash}X@{}}
\toprule
 & \textbf{Rep.\ agent $\to$ Het.\ agent} \\
\midrule
State & $k_t$ $\to$ $\lambda_t$ ($N$ bin masses) \\
Steady state & $k^{*}$ $\to$ $\bar{\lambda}$ ($N$ numbers) \\
Expansion & $\hat{k}_t$ $\to$ $\hat{\lambda}_t$ \\
\bottomrule
\end{tabularx}
\end{center}
\medskip
\begin{examplebox}\footnotesize
Notation: the literature often writes $\mu$ for the distribution. This course writes $\lambda$ for the invariant distribution and $Q$ for the operator that moves it (2.A3).
\end{examplebox}
\end{column}
\end{columns}
\vspace{1mm}
\src{Quant\_Macro Lec.~9, ``From Point to Distribution: The Conceptual Extension''.}
\end{frame}

%=============================================================================
\begin{frame}{Perturbing Around $\bar{\lambda}$: Four Steps}
\begin{columns}[T]
\begin{column}{0.55\textwidth}
\small
\begin{enumerate}\small\setlength{\itemsep}{1mm}
  \item \textbf{Compute the stationary distribution} $\bar{\lambda}=Q\bar{\lambda}$ --- exactly the fixed point of 2.A3, solved by iterating the Kolmogorov forward equation.
  \item \textbf{Discretize it}: a histogram with $N$ bins, so the state becomes the finite vector $(\lambda^1,\dots,\lambda^N,z)$.
  \item \textbf{Linearize} the law of motion for $\lambda_t$, and the household policies around steady-state prices.
  \item \textbf{Stack and apply Blanchard--Kahn}, exactly as in 5.A4. The output is first-order impulse responses for a heterogeneous-agent model.
\end{enumerate}
\medskip
The trade-off is a state vector of hundreds or thousands of bins --- which modern linear algebra handles, because the problem is still \emph{linear}.
\end{column}
\begin{column}{0.42\textwidth}
\begin{conceptbox}[title=\textbf{Why it works}]\small
If $F(\lambda_t,z_t)=0$ holds as an identity, we can differentiate it and match coefficients. The distribution simply has more components than a capital stock.
\end{conceptbox}
\medskip
\begin{examplebox}\footnotesize
\textbf{Reiter (2009)}; \textbf{Auclert, Bard\'oczy, Rognlie \& Straub (2021)} for the sequence-space Jacobian, which is the same idea organised around impulse responses instead of a state vector.
\end{examplebox}
\end{column}
\end{columns}
\vspace{1mm}
\src{Quant\_Macro Lec.~9, ``How to Perturb Around a Distribution'' and ``Intuition: Why This Works''.}
\end{frame}

%=============================================================================
\begin{frame}{Takeaways}
\begin{takeawaybox}
\begin{itemize}\small
  \item Second order is \textbf{12 equations in 12 unknowns}, \textbf{linear} given the first-order coefficients; every order above the first is bookkeeping. And every \textbf{odd} derivative in $\chi$ vanishes, so risk enters only through variances.
  \item The surviving object is the \textbf{risk correction} $\tfrac12 c_{\chi\chi}$: a constant shift that kills certainty equivalence and moves the ergodic distribution. Measured at $-0.00041\%$ of steady-state consumption, scaling like $\sigma^{2}$ to within $2\%$ at five times the shock --- and not at twenty.
  \item Second-order \textbf{simulations} feed squares back into squares and manufacture orders the solution never had. \textbf{Prune} (Kim--Kim--Schaumburg--Sims 2008).
  \item Measured: at this calibration pruning changes \textbf{nothing}, and at $30\times$ the shock the unpruned path explodes at $t=3{,}967$. Prune anyway --- the failure is silent until it is not.
  \item Validity is set by the \textbf{nearest singularity in the complex plane}. You cannot see it, so measure residuals off the expansion point instead.
  \item The expansion point can be a \textbf{distribution} --- Reiter (2009), Auclert et al.\ (2021).
\end{itemize}
\end{takeawaybox}
\end{frame}

%=============================================================================
\begin{frame}{Bridge: Now Put a Number on It}
\begin{columns}[T]
\begin{column}{0.53\textwidth}
\small
Block A has claimed a great deal: fast, accurate near the steady state, useless for risk at first order, valid over an unknown range. Every one of those is a measurable statement about the quarterly model we already solved globally in 4.B5.
\medskip

The next deck measures all of them --- and finds the point where perturbation stops being good enough, in capital units.
\end{column}
\begin{column}{0.44\textwidth}
\begin{conceptbox}[title=\textbf{Next (5.A6)}]\small
The worked example: the same model as 3.A4 and 4.B5, solved by perturbation in five steps, graded against the Session 4 global solution near and far from $k^{*}$.
\end{conceptbox}
\medskip
\begin{examplebox}\footnotesize
And the question Block B answers: what do you do when ``far from the steady state'' is where the economics is?
\end{examplebox}
\end{column}
\end{columns}
\end{frame}

\end{document}
